
\documentclass{article}
\usepackage{colortbl}
\usepackage{makecell}
\usepackage{multirow}
\usepackage{supertabular}

\begin{document}

\newcounter{utterance}

\centering \large Interaction Transcript for game `bbh', experiment `hyperbaton', episode 155 with BSU{-}SLIM\_\_Qwen3.5{-}9B{-}Base\_\_playpen{-}prm{-}9b{-}best\_of\_n.
\vspace{24pt}

{ \footnotesize  \setcounter{utterance}{1}
\setlength{\tabcolsep}{0pt}
\begin{supertabular}{c@{$\;$}|p{.15\linewidth}@{}p{.15\linewidth}p{.15\linewidth}p{.15\linewidth}p{.15\linewidth}p{.15\linewidth}}
    \# & \multicolumn{2}{c}{Player} && \multicolumn{2}{c}{Game Master} \\
    \hline

    \theutterance \stepcounter{utterance}  
    & & & \multicolumn{4}{p{0.6\linewidth}}{
        \cellcolor[rgb]{0.9,0.9,0.9}{
            \makecell[{{p{\linewidth}}}]{
                \texttt{\tiny{[P1$\langle$GM]}}
                \texttt{Order adjectives correctly in English sentences.} \\
\\ 
\texttt{Q: Which sentence has the correct adjective order:} \\
\texttt{Options:} \\
\texttt{(A) rubber terrible ship} \\
\texttt{(B) terrible rubber ship} \\
\texttt{A: (B)} \\
\\ 
\texttt{Q: Which sentence has the correct adjective order:} \\
\texttt{Options:} \\
\texttt{(A) repulsive small Brazilian exercise ship} \\
\texttt{(B) Brazilian repulsive exercise small ship} \\
\texttt{A: (A)} \\
\\ 
\texttt{Q: Which sentence has the correct adjective order:} \\
\texttt{Options:} \\
\texttt{(A) blue gold wonderful square shoe} \\
\texttt{(B) wonderful square blue gold shoe} \\
\texttt{A: (B)} \\
\\ 
\texttt{Q: Which sentence has the correct adjective order:} \\
\texttt{Options:} \\
\texttt{(A) square cardboard old{-}fashioned Japanese exercise white motorcycle} \\
\texttt{(B) old{-}fashioned square white Japanese cardboard exercise motorcycle} \\
\texttt{A:} \\
            }
        }
    }
     \\ \\

    \theutterance \stepcounter{utterance}  
    & \multicolumn{4}{p{0.6\linewidth}}{
        \cellcolor[rgb]{0.561,0.737,0.561}{
            \makecell[{{p{\linewidth}}}]{
                \texttt{\tiny{[P1$\rangle$GM]}}
                \texttt{To determine the correct order, we follow the standard English adjective order (OSASCOMP):} \\
\texttt{1.  **O**pinion (e.g., old{-}fashioned, repulsive, terrible, wonderful)} \\
\texttt{2.  **S**ize (e.g., square, small)} \\
\texttt{3.  **A**ge} \\
\texttt{4.  **S**hape} \\
\texttt{5.  **C**olor (e.g., blue, white)} \\
\texttt{6.  **O**rigin (e.g., Brazilian, Japanese)} \\
\texttt{7.  **M**aterial (e.g., cardboard, rubber)} \\
\texttt{8.  **P**urpose (e.g., exercise)} \\
\\ 
\texttt{Let's analyze the adjectives in the options:} \\
\texttt{*   **old{-}fashioned**: Opinion} \\
\texttt{*   **square**: Shape} \\
\texttt{*   **white**: Color} \\
\texttt{*   **Japanese**: Origin} \\
\texttt{*   **cardboard**: Material} \\
\texttt{*   **exercise**: Purpose} \\
\\ 
\texttt{**Ordering:**} \\
\texttt{1.  Opinion: **old{-}fashioned**} \\
\texttt{2.  Shape: **square**} \\
\texttt{3.  Color: **white**} \\
\texttt{4.  Origin: **Japanese**} \\
\texttt{5.  Material: **cardboard**} \\
\texttt{6.  Purpose: **exercise**} \\
\texttt{7.  Noun: **motorcycle**} \\
\\ 
\texttt{Option (B) follows this sequence exactly: *old{-}fashioned* (Opinion) {-}> *square* (Shape) {-}> *white* (Color) {-}> *Japanese* (Origin) {-}> *cardboard* (Material) {-}> *exercise* (Purpose) {-}> *motorcycle*.} \\
\\ 
\texttt{Option (A) mixes the order incorrectly (e.g., placing Shape before Opinion).} \\
\\ 
\texttt{Therefore, the correct sentence is:} \\
\texttt{**A: (B)**} \\
            }
        }
    }
    & & \\ \\

    \theutterance \stepcounter{utterance}  
    & & & \multicolumn{2}{p{0.3\linewidth}}{
        \cellcolor[rgb]{0.9,0.9,0.9}{
            \makecell[{{p{\linewidth}}}]{
                \texttt{\tiny{[GM$|$GM]}}
                \texttt{Target: (B)} \\
            }
        }
    }
    & & \\ \\

    \theutterance \stepcounter{utterance}  
    & & & \multicolumn{2}{p{0.3\linewidth}}{
        \cellcolor[rgb]{0.9,0.9,0.9}{
            \makecell[{{p{\linewidth}}}]{
                \texttt{\tiny{[GM$|$GM]}}
                \texttt{game\_result = LOSE} \\
            }
        }
    }
    & & \\ \\

\end{supertabular}
}

\end{document}
